TODO: move this

\begin{theorem}[Stability from coercive $\Phi$-regularization]
	\label{thm:stability-coercive-phi}
	Let $(X,\|\cdot\|_X)$ be a reflexive Banach space of functions on a domain $M$, and let
	\[
		R_\theta(f) := \langle \theta, \Phi(f) \rangle \in [0,\infty]
	\]
	where $\Phi:X\to V$ is a measurable map into a real locally convex space $V$ and $\theta\in V^\ast$ is such that $R_\theta$ is a proper, convex, lower semicontinuous functional on $X$.
	Assume the following:

	\begin{enumerate}
		\item \textbf{Coercivity (modulo a nullspace).}
		      There exists a finite-dimensional subspace $\mathcal N\subset X$ and a constant $c>0$ such that for all $f\in X$,
		      \[
			      R_\theta(f) \;\ge\; c\,\|f-\Pi_{\mathcal N}f\|_X - c,
		      \]
		      where $\Pi_{\mathcal N}$ denotes a fixed linear projection onto $\mathcal N$.
		      In particular, the sublevel sets $\{f: R_\theta(f)\le C\}$ are relatively compact in $X/\mathcal N$.

		\item \textbf{Strong convexity of the empirical loss.}
		      For each dataset $D=\{(x_i,y_i)\}_{i=1}^n$, define
		      \[
			      \mathcal L_D(f):=\frac{1}{n}\sum_{i=1}^n \ell\big(f(x_i),y_i\big),
		      \]
		      and assume $\mathcal L_D$ is $\alpha$-strongly convex on $X/\mathcal N$ with respect to $\|\cdot\|_X$, uniformly over all datasets $D$ (for some $\alpha>0$).

		\item \textbf{Lipschitz dependence on function values.}
		      Assume $\ell(\cdot,y)$ is $L$-Lipschitz in its first argument for all $y$, and that point-evaluation is continuous on $X$ in the sense that there exists $\kappa>0$ such that for all $x\in M$ and all $f,g\in X$,
		      \[
			      |f(x)-g(x)| \le \kappa \|f-g\|_X.
		      \]
	\end{enumerate}

	% Fix $\lambda>0$, and define the (regularized ERM / MAP) estimator
	% \[
	% \hat f_D \in \arg\min_{f\in X}\Big\{ \mathcal L_D(f) + \lambda R_\theta(f)\Big\},
	% \]
	% where $\hat f_D$ is chosen uniquely in $X/\mathcal N$ (e.g.\ by imposing a gauge condition such as $\Pi_{\mathcal N}\hat f_D=0$).

	% Then the estimator is \textbf{stable to masking} in the following sense:
	% if $D^{(j)}$ denotes the dataset obtained by removing (masking) the $j$-th sample from $D$, then for all $j\in\{1,\dots,n\}$,
	% \[
	% \|\hat f_D - \hat f_{D^{(j)}}\|_X
	% \;\le\;
	% \frac{2L\kappa}{\alpha\, n}.
	% \]
	% Consequently, the change in loss on any test point $(x,y)$ is uniformly bounded by
	% \[
	% \big|\ell(\hat f_D(x),y)-\ell(\hat f_{D^{(j)}}(x),y)\big|
	% \;\le\;
	% \frac{2L^2\kappa^2}{\alpha\, n}.
	% \]
	Observed data $y_i = f^*(x_i) + \mathcal{N}(0, \sigma^2)$
	\begin{align*}
		\hat f = \argmin_f \frac{1}{n}\sum_{i=1}^n \ell(f(x_i), y_i) + \lambda R(f) \Rightarrow
		\|f^* - \hat f\|_{X} \leq C(n, \sigma, R, \lambda)
	\end{align*}

\end{theorem}
\begin{proposition}[Bilevel learning of $\theta$ as a MAP surrogate]
	Let $R_\theta(f)=\langle \theta,\Phi(f)\rangle$ and define the inner estimator
	\[
		\hat f_\theta(D)\in\arg\min_{f\in X}\Big\{\mathcal L_D(f)+\lambda R_\theta(f)\Big\}.
	\]
	Consider the bilevel objective based on masking (or train/validation splits),
	\[
		\min_\theta \; \mathbb E_{\text{mask}}\Big[\mathcal L_{\mathrm{val}}\big(\hat f_\theta(D_{\mathrm{train}})\big)\Big]+\Omega(\theta).
	\]
	Then bilevel optimization selects $\theta$ so as to minimize the predictive risk of the MAP estimator induced by the Gibbs prior
	\[
		p_\theta(df)\propto \exp\!\big(-\lambda R_\theta(f)\big)\,\nu(df).
	\]
	In contrast, maximum-likelihood learning of $\theta$ for the induced Gibbs model includes an additional global complexity term $\log Z(\theta)$, which is generally not captured by the bilevel objective unless explicitly approximated (e.g., via Laplace/Occam corrections).
\end{proposition}



\begin{proposition}[Bilevel learning of $\theta$ as a MAP surrogate]
	Let $R_\theta(f)=\langle \theta,\Phi(f)\rangle$ and define the inner estimator
	\[
		\hat f_\theta(D)\in\arg\min_{f\in X}\Big\{\mathcal L_D(f)+\lambda R_\theta(f)\Big\}.
	\]
	Consider the bilevel objective based on masking (or train/validation splits),
	\[
		\min_\theta \; \mathbb E_{\text{mask}}\Big[\mathcal L_{\mathrm{val}}\big(\hat f_\theta(D_{\mathrm{train}})\big)\Big]+\Omega(\theta).
	\]
	Then bilevel optimization selects $\theta$ so as to minimize the predictive risk of the MAP estimator induced by the Gibbs prior
	\[
		p_\theta(df)\propto \exp\!\big(-\lambda R_\theta(f)\big)\,\nu(df).
	\]
	In contrast, maximum-likelihood learning of $\theta$ for the induced Gibbs model includes an additional global complexity term $\log Z(\theta)$, which is generally not captured by the bilevel objective unless explicitly approximated (e.g., via Laplace/Occam corrections).
\end{proposition}
